% \iffalse meta-comment
%
% hit-spread.dtx 是字间距撑开
%
% \fi
%
% \section{字间距撑开}
%
% \changes{v3.2a}{2026/08/07}{加 \cs{hit@spread@by:nn}，把标题的字间距从常量取值里
%   抽出来。原先撑开是写进取值的，同一个值还要供目录条目与 PDF 书签用，
%   那两处不该有这个空隙}
% 规范里「结论」「致谢」「目录」这类两字标题要撑开排。原先撑开是直接写进常量取值的：
% \begin{latex}
% conclusion-heading-zh = {结\hspace{\ccwd}论}
% \end{latex}
% 这样一个值供三处用，只有标题那处想要这个空隙。目录条目与 PDF 书签拿到的也是带空隙
% 的版本，书签里读出来是「结 论」。
%
% \cs{hit@spread@by:nn}\marg{字间距}\marg{文本} 在文本的每两个字之间插一段
% \cs{hspace}，与手写的写法逐个记号相同。常量里只留内容，撑开由排标题的地方做：
% \begin{latex}
% \hit@spread@by:nn { \ccwd } { \hit@conclusion@heading@zh }
% \end{latex}
%
% 参数二可以是文本，也可以是存着文本的宏，展开一次取到内容。只展开一次是有意的：
% 取值里若有 \cs{kaishu} 这类字体命令，全展开会出事。
%
% 参数一展开后是空的就原样输出，一个记号都不插。撑开多少是按变体定的
% （\option{heading-spread}），博士后那一档三个标题都不撑开，取的就是空值。
%
% \option{heading-spread} 的默认值是 \texttt{1em}，也就是一个全角字宽：Word 版
% 模板里「结论」两字之间是两个半角空格，宋体的半角空格是半个字宽，两个正好一个字。
% 不用 \cs{ccwd}——那是「字宽加字距」，\cs{frontmatter} 里的
% \cs{ziju}\marg{0.05580808} 一开，它就比一个字宽多 5.6\%。
%
%    \begin{macrocode}
%<*shared-spread>
%    \end{macrocode}
%
%    \begin{macrocode}
\tl_new:N \l__hit_spread_tl
\bool_new:N \l__hit_spread_first_bool
\cs_new_protected:Npn \hit@spread@by:nn #1#2
  {
    \tl_if_blank:oTF {#1} { #2 } { \__hit_spread_by:nn {#1} {#2} }
  }
\cs_new_protected:Npn \__hit_spread_by:nn #1#2
  {
    \tl_set:No \l__hit_spread_tl {#2}
    \bool_set_true:N \l__hit_spread_first_bool
    \tl_map_inline:Nn \l__hit_spread_tl
      {
        \bool_if:NTF \l__hit_spread_first_bool
          { \bool_set_false:N \l__hit_spread_first_bool }
          { \hspace {#1} }
        ##1
      }
  }
%    \end{macrocode}
%
% \changes{v3.2a}{2026/08/07}{加 \cs{hit@spread@to:nn}，封面字段标签改成分散对齐到
%   固定宽度，不再手算字间距}
% \cs{hit@spread@to:nn}\marg{宽度}\marg{文本} 把文本分散对齐到固定宽度。封面上
% 成组的字段标签要对齐到同一列宽，规范给的是那个列宽，不是字间距：
% \begin{latex}
% graduate-supervisor-field-label = {\hit@spread@to:nn{6em}{导师}}
% \end{latex}
% 原先写的是算完之后的字间距，四个字的那几个还要写成 \texttt{0.6666666} 倍，
% 改一次列宽得把同组的每一条重算一遍。
%
% \textbf{列宽用 \texttt{em} 不用 \cs{ccwd}。} 一个汉字的宽度是 \texttt{em}；
% \cs{ccwd} 是「字宽加字距」，开了 \option{zijv} 就比字宽大一截（本模板的
% \cs{frontmatter} 里有 \cs{ziju}\marg{0.05580808}，那之后 \cs{ccwd} 是字宽的
% 1.0558 倍）。标签里的字之间隔着别的东西，不走字距，所以按字宽算。
%
% 混用这两个单位正是原先那批取值出错的地方：\texttt{博\dots 姓名} 写成
% 「5 个字 + 4 段 0.5\cs{ccwd}」，\texttt{合作导师} 写成「4 个字 + 3 段 \cs{ccwd}」，
% 字算的是字宽、间距算的是 \cs{ccwd}，两条式子把同一个差额放大成不同的量，
% 一个 99.94bp、一个 100.72bp，值那一列对不齐。
%
% \changes{v3.2a}{2026/08/12}{分散对齐的三套实现并成一套，底层统一走
%   \pkg{xeCJKfntef} 的 \env{CJKfilltwosides}}
% \textbf{底层是 \env{CJKfilltwosides}。} 原先这里自己写了一遍：\cs{hbox\_to\_wd:nn}
% 套 \cs{tl\_map\_inline:Nn}，每两个记号之间塞一个 \cs{hfill}。那个写法在纯汉字上
% 与现在等价（45 变体逐字节相同），碰上标点就不对：\texttt{学院（部} 会被排成
% 「学\ 院\ （\ 部」，左括号后面多出一格。\env{CJKfilltwosides} 撑开的是 xeCJK 的
% \cs{CJKglue}，标点前后本来就没有这段胶，所以不会。
%
% 同一个道理，封面上原先还有第三套写法 \cs{makebox}\oarg{宽度}\oarg{s}，也并到
% 这里了。它撑开的同样是 \cs{CJKglue}，结果一致，只是内容撑不满时会报
% underfull box。
%    \begin{macrocode}
\cs_new_protected:Npn \hit@spread@to:nn #1#2
  { \begin { CJKfilltwosides } [b] {#1} #2 \end { CJKfilltwosides } }
%    \end{macrocode}
%
% \env{CJKfilltwosides} 是个 \env{minipage}，内容比给的宽度宽就折行。要的是「压进
% 这个宽度、不折行」的话它做不到，那种情况仍写 \cs{makebox}\oarg{宽度}\oarg{s}。
% \file{src/setup/hit-defaults.dtx} 里 \option{graduate-major-field-label} 的
% \texttt{学科/专业学位类别} 是唯一一处：九个字压进八个字宽。
%
% \changes{v3.2a}{2026/08/08}{两个宏各加一个不带冒号的名字，给 expl3 之外的取值用}
% 配置层的取值是在 \pkg{expl3} 之外记号化的（见 CONTRIBUTING §4.7），那里 \texttt{:}
% 不是字母，\cs{hit@spread@to:nn} 会被拆成 \cs{hit@spread@to} 加 \texttt{:nn}，
% 报 Undefined control sequence。所以两个宏各留一个不带冒号的名字，取值里写这个：
% \begin{latex}
% graduate-supervisor-field-label = {\hit@spread@to{6em}{导师}}
% \end{latex}
% 与 \cs{hit@if@option@TF} 同一个路数——凡是要在排版内容里出现的，名字里都不带冒号。
%    \begin{macrocode}
\cs_new_eq:NN \hit@spread@to \hit@spread@to:nn
\cs_new_eq:NN \hit@spread@by \hit@spread@by:nn
%    \end{macrocode}
%
%    \begin{macrocode}
%</shared-spread>
%    \end{macrocode}
%
% \Finale
